\chapter{Rozwiązania silne i słabe oraz twierdzenie Yamada--Watanabego}
\label{ch:weak-strong-sde}

\section{Motywacja: co należy ustalić przed rozwiązaniem równania?}

Gdy rozwiązujemy klasyczne równanie różniczkowe, znamy funkcję
wymuszającą i szukamy odpowiedniej trajektorii. W równaniu
stochastycznym sytuacja jest subtelniejsza. Możemy albo z góry
ustalić przestrzeń probabilistyczną, filtrację i proces Wienera,
a następnie szukać procesu $X$ na tej przestrzeni, albo
pozwolić, by przestrzeń i szum były częścią konstruowanego
rozwiązania. Pierwszy problem nazywamy \emph{silnym}, a drugi
\emph{słabym}. Przymiotniki te nie dotyczą rozmiaru współczynników
ani siły zastosowanej normy.

Co więcej, pytanie o \emph{istnienie} rozwiązania i pytanie
o jego \emph{jednoznaczność} są niezależne. Musimy również
odróżnić dwie jednoznaczności: zgodność trajektorii przy
wspólnym szumie oraz zgodność samych praw rozkładu. To
rozróżnienie ma znaczenie nawet wtedy, gdy rozważamy wyłącznie
równania o ciągłych trajektoriach.

\section{Dokładne definicje i podstawowe implikacje}

Niech $T>0$, $d,m\in\N$, $x\in\R^d$,
$b:[0,T]\times\R^d\to\R^d$ oraz
$\sigma:[0,T]\times\R^d\to\R^{d\times m}$ będą
funkcjami borelowskimi. W tym rozdziale rozważamy
\[
 X_t=x+\int_0^t b(s,X_s)\,\dd s
 +\int_0^t\sigma(s,X_s)\,\dd W_s,
 \qquad 0\le t\le T.
 \tag{14.1}\label{eq:weak-strong-main}
\]

\begin{definicja}[Rozwiązanie słabe]
\label{def:weak-solution}
\emph{Rozwiązaniem słabym} równania \eqref{eq:weak-strong-main}
nazywamy układ
$(\Omega,\calF,\Prob,(\calF_t)_{0\le t\le T},W,X)$,
w którym $(\Omega,\calF,\Prob)$ jest przestrzenią
probabilistyczną, filtracja spełnia warunki zwykłe,
$W$ jest $m$-wymiarowym procesem Wienera \emph{względem tej
filtracji}, a $X$ jest ciągłym, adaptowanym procesem
$\R^d$-wartościowym o $X_0=x$ i
\[
 \int_0^T\bigl(\|b(s,X_s)\|+
 \|\sigma(s,X_s)\|_{\mathrm{HS}}^2\bigr)\,\dd s<\infty
 \quad\text{prawie na pewno}.
\]
Wymagamy, aby \eqref{eq:weak-strong-main} zachodziło
prawie na pewno dla wszystkich $t$ jednocześnie.
Przestrzeń probabilistyczna i proces Wienera są częścią
poszukiwanego rozwiązania.
\end{definicja}

\begin{definicja}[Rozwiązanie silne]
\label{def:strong-solution-precise}
\emph{Silnym rozwiązaniem} na zadanej przestrzeni z zadanym
procesem Wienera $W$ nazywamy rozwiązanie
\eqref{eq:weak-strong-main}, adaptowane do uzupełnionej,
prawostronnie ciągłej filtracji naturalnej procesu $W$.
Równoważnie, w powyższym przypadku deterministycznego stanu
początkowego, można wybrać proces $X$ będący nieantycypującym,
mierzalnym funkcjonałem trajektorii $W$. Dla losowego stanu
początkowego należy do danych tego funkcjonału dodać stan
początkowy, niezależny od przyszłych przyrostów szumu.
\end{definicja}

\begin{definicja}[Dwa rodzaje jednoznaczności]
\label{def:path-law-uniq}
Równanie \eqref{eq:weak-strong-main} ma
\emph{jednoznaczność ścieżkową}, jeżeli na każdej wspólnej
przestrzeni probabilistycznej z tym samym procesem Wienera
każde dwa rozwiązania $X,Y$ o tym samym stanie początkowym
są nierozróżnialne, to jest
$\Prob(X_t=Y_t\text{ dla wszystkich }t\le T)=1$.

Mówimy o \emph{jednoznaczności co do prawa}, jeżeli prawa
procesu $X$ na $C([0,T];\R^d)$ są jednakowe dla wszystkich
słabych rozwiązań tego równania o wskazanym stanie początkowym.
Nie jest to z definicji to samo, co jednoznaczność wspólnego
prawa pary $(X,W)$.
\end{definicja}

\begin{uwaga}[Nie mylimy przestrzeni praw z przestrzenią trajektorii]
Równość $\Law(X)=\Law(Y)$ jest równością miar na
$C([0,T];\R^d)$ i nie oznacza, że $X_t(\omega)=Y_t(\omega)$.
Nawet jeśli dwa procesy mają to samo prawo, mogą być
zdefiniowane na różnych przestrzeniach probabilistycznych.
Jednoznaczność ścieżkowa wymaga porównania dwóch procesów
na \emph{tej samej} przestrzeni i przy \emph{tym samym} szumie.
\end{uwaga}

\begin{przyklad}[Istnienie bez jednoznaczności]
Niech $d=m=1$, $\sigma\equiv0$,
$b(x)=2\sqrt{\max\{x,0\}}$ i $x=0$.
Dla każdego $c\ge0$ funkcja
$X_t^{(c)}=(t-c)_+^2$ spełnia
$X_t^{(c)}=\int_0^t2\sqrt{X_s^{(c)}}\,\dd s$.
Istnieje też rozwiązanie identycznie zerowe.
Są to rozwiązania silne na każdej przestrzeni z procesem
Wienera, ale nie zachodzi ani jednoznaczność ścieżkowa,
ani jednoznaczność co do prawa. Nie można więc zastąpić
hipotezy Lipschitza samą ciągłością współczynnika.
\end{przyklad}

\section{Przygotowanie do twierdzenia Yamada--Watanabego}

Dowód nie polega na formalnej manipulacji z prawami
rozkładu. Kluczowa idea jest następująca: \emph{dwa
warunkowo niezależne rozwiązania przy tej samej realizacji
szumu muszą, na mocy jednoznaczności ścieżkowej, być równe}.
Z tego wywnioskujemy, że przy ustalonym szumie rozkład
warunkowy rozwiązania skupia się w jednym punkcie.
Aby uniknąć ukrytego założenia o mierzalności, osobno
uzasadnimy wszystkie kroki tej konstrukcji.

\begin{lemat}[Niezależność przyszłych przyrostów i przyczynowość]
\label{lem:YW-causality}
Niech $(W,X)$ będzie słabym rozwiązaniem
\eqref{eq:weak-strong-main}. Wybierzmy regularny rozkład
warunkowy $K(w,\dd y)$ całej trajektorii $X$ pod warunkiem
$W=w$, gdzie $w\in C([0,T];\R^m)$ i
$y\in C([0,T];\R^d)$. Dla każdego $t\in[0,T]$ oraz
każdej ograniczonej borelowskiej funkcji $F$ zależnej
wyłącznie od $y|_{[0,t]}$ istnieje wersja funkcji
$w\mapsto\int F(y)K(w,\dd y)$ mierzalna względem
$\sigma(w_s:0\le s\le t)$ po uzupełnieniu miary Wienera.
\end{lemat}
\begin{proof}
Zmienna $F(X)$ jest $\calF_t$-mierzalna.
Przyszły proces przyrostów $(W_{t+s}-W_t)_{0\le s\le T-t}$
jest niezależny od $\calF_t$, a więc także od $F(X)$
oraz historii $W$ do chwili $t$. Zatem dla każdej ograniczonej
funkcji $G$ przeszłości $W$ i każdej ograniczonej funkcji $H$
przyszłych przyrostów mamy
\[
 \E[F(X)G(W_{\le t})H(W_{t+\cdot}-W_t)]
 =\E[F(X)G(W_{\le t})]\,\E H(W_{t+\cdot}-W_t).
\]
Twierdzenie o klasie monotonicznej rozciąga tę równość
na ograniczone funkcje całej trajektorii $W$. Oznacza to,
że $\E[F(X)\mid\sigma(W_s:0\le s\le T)]$ jest równe
$\E[F(X)\mid\sigma(W_s:0\le s\le t)]$ prawie na pewno.
Lewa strona jest $\int F(y)K(W,\dd y)$, co daje tezę.
\end{proof}

\begin{lemat}[Równanie Itô jako własność wspólnego prawa]
\label{lem:YW-transfer}
Jeśli dwie pary ciągłych procesów $(X,W)$ i
$(\widehat X,\widehat W)$ mają to samo prawo na
$C([0,T];\R^d)\times C([0,T];\R^m)$, a oba procesy
$W,\widehat W$ są procesami Wienera względem odpowiednich
filtracji zawierających historie $X,\widehat X$, to
spełnianie równania \eqref{eq:weak-strong-main} przez
pierwszą parę implikuje jego spełnianie przez drugą
(z tymi samymi współczynnikami i stanem początkowym).
\end{lemat}
\begin{proof}
Całka $\int_0^t b(s,X_s)\,\dd s$ jest borelowskim
funkcjonałem trajektorii $X$ tam, gdzie istnieje.
Dla całki stochastycznej oznaczmy
$H_s=\sigma(s,X_s)$ i najpierw obetnijmy jej wartości,
$H_s^{(K)}=H_s\ind{\{\|H_s\|\le K\}}$.
Dla ustalonego $K$ uśrednienia lewostronne
\[
 H_s^{(K,n)}=n\int_{(s-1/n)^+}^{s}H_r^{(K)}\,\dd r
\]
są ciągłymi, adaptowanymi funkcjami czasu i zbieżne do
$H^{(K)}$ w $L^2(\Prob\otimes\dd s)$ na podstawie
ciągłości przesunięć w $L^2$ oraz zbieżności zdominowanej.
W dalszym przybliżeniu zastępujemy je wartościami na
lewych końcach drobnego deterministycznego podziału.
Otrzymana całka jest skończoną sumą
$\sum_j H_{t_j}^{(K,n)}(W_{t_{j+1}}-W_{t_j})$,
a więc \emph{borelowskim funkcjonałem pary} $(X,W)$.
Izometria Itô daje zbieżność tych sum do
$\int_0^tH_s^{(K)}\,\dd W_s$ w prawdopodobieństwie.
Po dopuszczeniu $K\to\infty$ i lokalizacji przy
$\int_0^T\|H_s\|^2\,\dd s$ otrzymujemy całą całkę.
Wybierając podciągi, wyrażamy zatem całkę w każdej
wymiernej chwili jako prawie wszędzie określoną granicę
borelowskich funkcjonałów pary $(X,W)$.
Ta sama konstrukcja i te same rozkłady par pozwalają
przenieść tożsamość równania na $(\widehat X,\widehat W)$.
Ciągłość obu stron rozszerza ją ze zbioru chwil wymiernych
na wszystkie $t$.
\end{proof}

\begin{twierdzenie}[Yamada--Watanabe: wersja dla równań Browna]
\label{th:YW}
Niech $b,\sigma$ będą funkcjami borelowskimi określonymi
przed \eqref{eq:weak-strong-main}, a $x\in\R^d$ będzie
ustalony. Załóżmy, że istnieje przynajmniej jedno słabe
rozwiązanie \eqref{eq:weak-strong-main}, i że zachodzi
jednoznaczność ścieżkowa w sensie
Definicji~\ref{def:path-law-uniq}. Wówczas:
\begin{enumerate}
\item istnieje silne rozwiązanie dla każdego $m$-wymiarowego
procesu Wienera na przestrzeni z filtracją naturalną,
po jej zwykłym uzupełnieniu;
\item rozwiązanie jest nieantycypującym mierzalnym funkcjonałem
trajektorii $W$, z dokładnością do nierozróżnialności;
\item zachodzi jednoznaczność co do prawa procesu $X$,
a nawet jednoznaczność wspólnego prawa pary $(X,W)$.
\end{enumerate}
\end{twierdzenie}
\begin{proof}
\emph{Krok 1: warunkowe kopie.} Ustalmy słabe rozwiązanie
i oznaczmy jego wspólne prawo $(W,X)$ przez $\mu$.
Ponieważ przestrzenie ciągłych trajektorii są polskie,
istnieje regularny rozkład warunkowy $K(w,\dd y)$
rozwiązania przy ustalonej trajektorii szumu.
Na przestrzeni kanonicznej trójek $(w,y_1,y_2)$
wprowadzamy miarę
\[
 \mathsf Q(\dd w,\dd y_1,\dd y_2)
 =\mu_W(\dd w)K(w,\dd y_1)K(w,\dd y_2),
\]
gdzie $\mu_W$ jest miarą Wienera. Warunkowo względem
$w$ procesy $y_1,y_2$ są niezależnymi kopiami z prawa
$K(w,\cdot)$.

\emph{Krok 2: wspólny proces Wienera.} Nie możemy po prostu
założyć, że $w$ jest procesem Wienera względem filtracji
zawierającej historie obu $y_i$: trzeba to udowodnić.
Niech $\calG_t$ będzie uzupełnioną filtracją generowaną
przez $w_s,y_{1,s},y_{2,s}$ dla $s\le t$ (następnie bierzemy
jej wersję prawostronnie ciągłą). Dla ograniczonej funkcji
$F$ zależnej od tych historii jej warunkowa średnia pod
warunkiem całego $w$ jest funkcją historii $w_{\le t}$,
ponieważ warunkowe prawo dwóch kopii jest iloczynem
$K(w,\cdot)$ i zachodzi Lemat~\ref{lem:YW-causality}.
Przyszłe przyrosty $w$ są niezależne od jego przeszłości,
zatem niezależne od $\calG_t$. Przejście do filtracji
prawostronnie ciągłej uzasadnia się przez przybliżenie
chwil $t$ z góry i prawostronną ciągłość trajektorii.
Proces $w$ pozostaje więc procesem Wienera względem
wspólnej filtracji.

\emph{Krok 3: obie kopie rozwiązują to samo równanie.}
Z definicji $\mathsf Q$ każda para $(y_i,w)$ ma prawo
$\mu$. Lematem~\ref{lem:YW-transfer} przenosimy zatem
równanie Itô na każdą z par, względem wspólnej filtracji
z kroku 2. Obie kopie mają wspólny szum $w$ i wspólny
stan początkowy. Z jednoznaczności ścieżkowej wynika
$y_1=y_2$ jako trajektorie, $\mathsf Q$-prawie na pewno.

\emph{Krok 4: skupienie rozkładu warunkowego w punkcie.}
Warunkowa niezależność i poprzednia równość dają
\[
 \int K(w,\dd y_1)K(w,\dd y_2)
 \ind{\{y_1\ne y_2\}}=0
 \quad\text{dla }\mu_W\text{-prawie każdego }w.
\]
Lemat~\ref{lem:dirac-audit}, udowodniony poniżej,
pokazuje, że taka miara musi być miarą Diraca.
Wobec tego $K(w,\cdot)=\delta_{F(w)}$ prawie wszędzie,
dla pewnego mierzalnego odwzorowania
$F:C([0,T];\R^m)\to C([0,T];\R^d)$.
Lemat~\ref{lem:YW-causality} pokazuje, że ograniczenie
$F(w)|_{[0,t]}$ zależy, z dokładnością do zbiorów zerowych,
wyłącznie od $w|_{[0,t]}$. Stąd proces $F(W)$ jest
adaptowany do uzupełnionej filtracji naturalnej i,
na mocy Lematu~\ref{lem:YW-transfer}, spełnia SDE.
Otrzymaliśmy silne istnienie. To samo odwzorowanie możemy
zastosować do dowolnego innego procesu Wienera, gdyż jego
prawo jest $\mu_W$.

\emph{Krok 5: jednoznaczność prawa.} Dla dwóch dowolnych
słabych rozwiązań wybierzmy ich jądra warunkowe
$K_1(w,\cdot)$ i $K_2(w,\cdot)$ i powtórzmy kroki 1--3,
tym razem losując jedną kopię z $K_1$, a drugą z $K_2$.
Jednoznaczność ścieżkowa znów wymusza równość kopii.
Ponieważ obie są warunkowo niezależne, ich jądra są
jednocześnie jednakowymi miarami Diraca. Otrzymujemy
to samo prawo pary $(X,W)$, a po marginalizacji ---
to samo prawo $X$.
\end{proof}

\begin{uwaga}[Zakres twierdzenia]
Powyższy dowód korzysta ze standardowego twierdzenia o
istnieniu regularnych rozkładów warunkowych na przestrzeniach
polskich oraz z konstrukcji miary na iloczynie za pomocą jąder.
Nie jest to bezpośredni wniosek z elementarnego rachunku
Itô; potrzebny jest też Lemat~\ref{lem:YW-transfer}.
Analogiczna idea działa dla wielu równań napędzanych
procesami ze skokami, ale wymaga ponownej kontroli
przyczynowości i kanonicznej przestrzeni losowych miar.
\end{uwaga}


\section{Doprecyzowanie argumentu warunkowego w twierdzeniu Yamada--Watanabego}

\begin{lemat}[Kiedy warunkowy rozkład jest miarą Diraca?]
\label{lem:dirac-audit}
Niech $S$ będzie przestrzenią polską, $K$ miarą
probabilistyczną na $S$ i niech $Y_1,Y_2$ będą
niezależnymi zmiennymi o rozkładzie $K$. Jeżeli
$\Prob(Y_1=Y_2)=1$, to $K=\delta_y$ dla pewnego $y\in S$.
Jeśli ponadto $w\mapsto K(w,\cdot)$ jest jądrem
probabilistycznym przyjmującym prawie wszędzie miary
Diraca, to jego atom $F(w)$ można wybrać mierzalnie
względem uzupełnienia odpowiedniej miary bazowej.
\end{lemat}
\begin{proof}
Wybieramy przeliczalną rodzinę zbiorów borelowskich
rozdzielającą punkty $S$ (na przykład kule o wymiernych
promieniach i środkach z przeliczalnego zbioru gęstego).
Dla każdego zbioru $A$ z tej rodziny niezależność daje
$\Prob(Y_1\in A,Y_2\notin A)=K(A)(1-K(A))=0$.
Wszystkie te wartości są więc równe zeru lub jedności.
Dwie niezależne realizacje należą prawie na pewno do
identycznych elementów całej rodziny, a rodzina rozdziela
punkty, więc są równe. Gdyby $K$ nie było miarą Diraca,
istniałby zbiór borelowski o prawdopodobieństwie ściśle
między zerem a jednością: można wziąć dwie rozłączne
okolice punktów należących do nośnika $K$.
Otrzymujemy sprzeczność.

Dla mierzalności rozważamy odwzorowanie
$\delta:S\to\mathcal P(S)$, $y\mapsto\delta_y$.
Jest ono ciągłe i różnowartościowe. Jego obraz jest
nawet domknięty: jeśli $\delta_{y_n}$ zbiega słabo do
$\mu$, to ciasność zbieżnej rodziny daje wspólny
zwarty zbiór zawierający prawie wszystkie $y_n$;
podciąg $y_{n_k}\to y$ pokazuje $\mu=\delta_y$.
Każdy zbieżny podciąg punktów ma tę samą granicę,
więc $y_n\to y$. Odwrotność obrazu jest zatem
ciągła, a w szczególności borelowska. Na zbiorze pełnej miary, na którym
$K(w,\cdot)$ jest miarą Diraca, kładziemy
$F(w)=\delta^{-1}(K(w,\cdot))$; poza nim wybieramy
ustalony punkt. To daje wymagany funkcjonał.
\end{proof}

\begin{uwaga}[Dlaczego sama równość praw nie wystarcza?]
W dowodzie Twierdzenia~\ref{th:YW} wspólne prawo pary
$(X,W)$ nie jest zastępnikiem warunku adaptacji.
Najpierw potrzebujemy Lematu~\ref{lem:YW-causality}, aby
udowodnić, że warunkowe kopie nie mają informacji o
przyszłych przyrostach $W$. Dopiero wtedy wspólny proces
$W$ pozostaje procesem Wienera względem filtracji
powiększonej o historie obu kopii i wolno zastosować
jednoznaczność ścieżkową. Lemat~\ref{lem:dirac-audit}
uzasadnia ostatni etap, w którym warunkowe prawo
zostaje zastąpione mierzalnym funkcjonałem $F(W)$.
\end{uwaga}

\section{Przykłady i zastosowania jednoznaczności}

\begin{zadanie}[Silne rozwiązanie jest automatycznie słabym]
Przypuśćmy, że współczynniki $b,\sigma$ spełniają
globalny warunek Lipschitza i liniowego wzrostu.
Co dokładnie wynika o wszystkich trzech rodzajach
jednoznaczności?

\emph{Rozwiązanie.} Z twierdzenia konstrukcyjnego
rozdziału~5 istnieje silne rozwiązanie i zachodzi
jednoznaczność ścieżkowa. W szczególności rozwiązanie
silne jest słabym, bo jego filtracja i szum spełniają
Definicję~\ref{def:weak-solution}. Hipotezy
Twierdzenia~\ref{th:YW} są spełnione, a więc zachodzi
jednoznaczność co do prawa, także dla dowolnych słabych
rozwiązań konstruowanych na innych przestrzeniach.
Rozwiązanie silne nie jest z definicji ``bardziej
regularną'' trajektorią; chodzi wyłącznie o źródło
losowości i mierzalność względem niego.
\end{zadanie}

\begin{twierdzenie}[Słabe istnienie przy ograniczonym dryfie]
\label{th:bounded-borel-weak}
Niech $b:[0,T]\times\R\to\R$ będzie ograniczoną funkcją
borelowską, a $x\in\R$. Równanie
\[
 X_t=x+\int_0^t b(s,X_s)\,\dd s+W_t
 \tag{14.2}\label{eq:borel-drift}
\]
ma słabe rozwiązanie, a jego prawo na $C([0,T];\R)$
jest jednoznaczne.
\end{twierdzenie}
\begin{proof}
\emph{Istnienie.} Na przestrzeni z procesem Wienera $B$
pod miarą $\mathsf P$ połóżmy $X_t=x+B_t$ i
\[
 Z_T=\exp\left(\int_0^Tb(s,X_s)\,\dd B_s
 -\tfrac12\int_0^Tb(s,X_s)^2\,\dd s\right).
\]
Ograniczoność $b$ zapewnia warunek Novikova,
więc $\E_{\mathsf P}Z_T=1$ i możemy określić
$\dd\mathsf Q=Z_T\dd\mathsf P$.
Twierdzenie Girsanowa z rozdziału~8 mówi, że
$W_t=B_t-\int_0^tb(s,X_s)\dd s$ jest
procesem Wienera pod $\mathsf Q$.
Wobec tego $X=x+\int b(s,X_s)\dd s+W$,
co daje słabe rozwiązanie.

\emph{Jednoznaczność prawa.} Niech $\widehat X$ będzie
dowolnym słabym rozwiązaniem napędzanym przez
$\widehat W$ pod miarą $\widehat{\mathsf P}$.
Girsanow zastosowany z przeciwnym znakiem do
$\theta_s=b(s,\widehat X_s)$ daje miarę
$\widehat{\mathsf Q}$, równoważną
$\widehat{\mathsf P}$, pod którą
$\widehat X_t-x=\widehat W_t+
\int_0^tb(s,\widehat X_s)\dd s$ jest
standardowym procesem Wienera. Jednocześnie
\[
 \frac{\dd\widehat{\mathsf P}}{\dd\widehat{\mathsf Q}}
 =\exp\left(\int_0^T b(s,\widehat X_s)\,\dd\widehat X_s
 -\frac12\int_0^Tb(s,\widehat X_s)^2\,\dd s\right).
\]
Całka po prawej jest całką Itô względem Browna
$\widehat X-x$ pod $\widehat{\mathsf Q}$ i jest
mierzalnym funkcjonałem trajektorii $\widehat X$
z dokładnością do zbioru zerowego, co uzasadnia
Lemat~\ref{lem:YW-transfer}. Ponieważ pod
$\widehat{\mathsf Q}$ prawo $\widehat X$ jest zawsze
prawem procesu $x+B$, powyższy wzór wyznacza jednoznacznie
jego prawo pod $\widehat{\mathsf P}$.
\end{proof}

\begin{uwaga}[Czego twierdzenie o ograniczonym dryfie nie dowodzi?]
W Twierdzeniu~\ref{th:bounded-borel-weak} skonstruowano
\emph{słabe} rozwiązanie i wykazano jednoznaczność co do
prawa. Nie udowodniono na tej podstawie jednoznaczności
ścieżkowej ani istnienia silnego rozwiązania przy dowolnym
ograniczonym borelowskim dryfie. Takie wzmocnienia wymagają
osobnych argumentów regularizacyjnych.
\end{uwaga}

\begin{przyklad}[Dlaczego sama jednoznaczność prawa nie zastępuje ścieżkowej?]
Dwa procesy $X,Y$ mogą być niezależnymi ruchami Browna.
Wówczas $\Law(X)=\Law(Y)$, ale
$\Prob(X_t=Y_t\text{ dla każdego }t\in[0,T])=0$.
Ta obserwacja dotyczy ogólnie praw procesów. Nie jest
kontrprzykładem do Twierdzenia~\ref{th:YW}, ponieważ
procesy nie zostały przedstawione jako dwa rozwiązania
tego samego SDE napędzanego przez \emph{wspólny} proces
Wienera.
\end{przyklad}

\begin{center}
\begin{tikzpicture}[>=Latex,node distance=2cm]
\node[draw,rounded corners,align=center,minimum width=3cm,minimum height=.7cm]
 (w) {zadany szum $W$};
\node[draw,rounded corners,align=center,below left=1.3cm and 1.9cm of w,
 minimum width=3.3cm,minimum height=.9cm] (strong)
 {rozwiązanie silne\\$X=F(W)$};
\node[draw,rounded corners,align=center,below right=1.3cm and 1.9cm of w,
 minimum width=3.4cm,minimum height=.9cm] (weak)
 {rozwiązanie słabe\\$\Law(W,X)$};
\draw[->,thick,blue!70!black] (w)--(strong)
 node[midway,left,scale=.85] {konstrukcja};
\draw[->,thick,blue!70!black] (w)--(weak)
 node[midway,right,scale=.85] {wybór prawa};
\draw[->,thick,dashed] (strong)--(weak)
 node[midway,below,scale=.85] {implikuje};
\end{tikzpicture}
\par\smallskip\small
Różnica między ujęciem silnym i słabym: w pierwszym
szum jest dany przed konstrukcją rozwiązania, a w drugim
możemy wybierać go razem z przestrzenią probabilistyczną.
\end{center}

\section{Podsumowanie i kierunki dalszej teorii}

Otrzymaliśmy wnikliwe rozróżnienie istnienia silnego i słabego,
jednoznaczności ścieżkowej oraz jednoznaczności co do prawa.
Twierdzenie Yamada--Watanabego wyjaśnia, dlaczego przy
istnieniu słabego rozwiązania jednoznaczność ścieżkowa
wymusza silne istnienie: rozkład warunkowy rozwiązania
przy danej trajektorii szumu musi się skoncentrować w jednym
punkcie. Dowód wymagał regularnych rozkładów warunkowych,
argumentu przyczynowości i kontroli całek Itô przy
przenoszeniu praw par procesów.

W kolejnych rozdziałach rozwiniemy całki po trajektoriach
w języku miary Wienera i wprowadzimy fizyczne całki
Feynmana, z uwzględnieniem różnicy pomiędzy miarą
probabilistyczną i oscylacyjnym formalizmem propagatora.




% ===========================================================================
% ROZDZIAŁ 15: całki po trajektoriach, most Browna, jądra Feynmana--Kaca
% ===========================================================================
% Część końcowa: ścisłe całki funkcyjne i rachunek operatorowy.
% W całym rozdziale 16 twierdzenie spektralne, klasyczne jądra gaussowskie
% i podstawowa teoria generatorów półgrup są jawnie przyjmowanymi narzędziami.
